\section{Detailed Task Description}
\begin{taskbox}{Problem 1 (Lasso Regularization Path)}
\small

Given a feature matrix
$X \in \mathbb{R}^{n \times p}$,
response $y \in \mathbb{R}^{n}$,
and a decreasing sequence
$\lambda_1 > \cdots > \lambda_K$,
define
\[
F_k(w)
:=
\frac{1}{2n}\|y-Xw\|_2^2
+
\lambda_k\|w\|_1,
\qquad
w_k^\star
\in
\arg\min_{w\in\mathbb{R}^p} F_k(w).
\]

A candidate solver returns approximate coefficients
$\widetilde{W}
=
(\widetilde{w}_1,\ldots,\widetilde{w}_K)$.
It passes the benchmark's objective-value check if
\[
F_k(\widetilde{w}_k)
\le
F_k(w_{k,\mathrm{sklearn}})
+
10^{-6},
\qquad
\text{for every } k,
\]
where correctness is checked on fresh instances distinct from the timing
instances. If any required check fails, the search score is zero.

Otherwise, letting $\mathcal{I}$ denote the timing instances and $t_i$ the
time required to compute the complete regularization path on instance $i$,
the search score is
\[
R_{\mathrm{search}}
=
\left(
\prod_{i\in\mathcal{I}} t_i
\right)^{-1/|\mathcal{I}|}.
\]

\end{taskbox}

\begin{taskbox}{Problem 2 (Sum--Difference Problem)}
\small

The Sum--Difference Problem asks for a finite set
$A \subset \mathbb{Z}$
whose normalized sumset is large relative to its normalized difference set.
The objective is
\[
\Gamma(A)
:=
\frac{
\log\!\left(|A+A|/|A|\right)
}{
\log\!\left(|A-A|/|A|\right)
},
\]
where
\[
A+A
:=
\{a+a' : a,a' \in A\},
\qquad
A-A
:=
\{a-a' : a,a' \in A\}.
\]

The goal is to construct a finite set $A \subset \mathbb{Z}$ that maximizes
$\Gamma(A)$.

\end{taskbox}


\begin{taskbox}{Problem 3 (Circle Packing in a Unit Square)}
\small

For $n \in \{26,32\}$, the circle-packing task asks for centers
$(x_i,y_i)\in[0,1]^2$
and radii $r_i\ge 0$ such that every circle lies inside the unit square
and no two circles overlap:
\[
r_i \le x_i \le 1-r_i,
\qquad
r_i \le y_i \le 1-r_i,
\]
\[
(x_i-x_j)^2+(y_i-y_j)^2
\ge
(r_i+r_j)^2,
\qquad
1\le i<j\le n.
\]

The objective is to maximize
\[
\sum_i r_i.
\]

\end{taskbox}

\begin{taskbox}{Problem 4 (Autocorrelation Inequalities)}
\small

For an integrable function $f:\mathbb{R}\to\mathbb{R}$, define the
autoconvolution
\[
(f*f)(t)
:=
\int_{\mathbb{R}} f(t-x)f(x)\,dx,
\qquad
t\in\left[-\frac12,\frac12\right].
\]

\textbf{First Autocorrelation Inequality.}
Find a non-negative integrable function
$f:\mathbb{R}\to\mathbb{R}$
supported on
$\left[-\frac14,\frac14\right]$
such that
\[
\int_{-1/4}^{1/4} f(x)\,dx = 1.
\]
The objective is to minimize
\[
\Phi_1(f)
:=
\max_{t\in[-1/2,\,1/2]} (f*f)(t).
\]

\textbf{Second Autocorrelation Inequality.}
Find a non-negative integrable function
$f:\mathbb{R}\to\mathbb{R}$
supported on
$\left[-\frac14,\frac14\right]$
such that
\[
\int_{-1/4}^{1/4} f(x)\,dx = 1.
\]
The objective is to maximize
\[
\Phi_2(f)
:=
\frac{\|f*f\|_2^2}
{\|f*f\|_1\|f*f\|_\infty}.
\]

\textbf{Third Autocorrelation Inequality.}
Find an integrable (possibly signed) function
$f:\mathbb{R}\to\mathbb{R}$
supported on
$\left[-\frac14,\frac14\right]$
such that
\[
\int_{-1/4}^{1/4} f(x)\,dx = 1.
\]
The objective is to minimize
\[
\Phi_3(f)
:=
\max_{t\in[-1/2,\,1/2]}
\left|(f*f)(t)\right|.
\]

\end{taskbox}



\section{Prompts}
\label{app:prompts}

For reproducibility, we provide the prompts used for online exploration and
replay-based exploration-policy improvement.
Variables enclosed by dollar signs or braces are instantiated by the runtime
system before execution.

\subsection{Exploration Prompt}
\label{app:exploration-prompt}

The following prompt is used to guide the discovery agent during online
exploration. It requires the agent to inspect the complete available discovery
history before proposing a new solution, explicitly reason about both successful
and failed attempts, and avoid repeatedly exploiting a locally saturated
direction.

\lstinputlisting[
    style=promptstyle,
    caption={Prompt used for online exploration.},
    label={lst:exploration-prompt}
]{appendix/prompt/exploration_prompt.txt}


\subsection{Replay-Based Policy Improvement Prompt}
\label{app:replay-prompt}

The following prompt is used by the controller-development agent during
historical replay. The agent modifies the exploration policy using feedback
obtained from replay over previously collected discovery trajectories while
remaining restricted to prefix-observable information.

\lstinputlisting[
    style=promptstyle,
    caption={Prompt used for replay-based improvement of the exploration policy.},
    label={lst:replay-prompt}
]{appendix/prompt/replay_learning_prompt.txt}

\section{Discovered Programs}
\label{app:discovered-lasso}

We provide the complete implementation of the Lasso-path solver discovered by
\textsc{Dream-RSI}. As discussed in Section~\ref{subsec:lasso}, the solver combines
strong-rule screening with adaptive Cauchy--Schwarz KKT pruning, disjoint
active-set bookkeeping, lazy Gram-matrix construction, and hardware-aware
optimizations.

\lstinputlisting[
    style=appendixcode,
    language=Python,
    caption={Complete Lasso-path solver discovered by \textsc{Dream-RSI}.},
    label={lst:discovered-lasso}
]{appendix/discovered_programs/lasso/code.cpp}